% ============================================================
%  TERM PROJECT – EXAMPLE DOCUMENTATION
%  Databases · THGA Bochum
%
%  This is a worked example of the *documentation* deliverable.
%  It shows the expected structure and level of detail. Your own
%  documentation describes your own application.
% ============================================================
\documentclass[11pt,a4paper]{article}

\usepackage{thga-db}

\renewcommand{\thgaKurs}{Databases}
\renewcommand{\thgaDatum}{Summer Term 2026}

\begin{document}
\thispagestyle{empty}
\thgaTitel{Term Project – Documentation (Example)}%
{Club Library: a database-backed lending system}

\begin{infobox}
\textbf{Author:} Jane Doe (example) \\
\textbf{Module:} Introduction to Database Management Systems · THGA Bochum \\
\textbf{Stack:} PostgreSQL · FastAPI · Docker Compose · tkinter (\texttt{.deb}) \\
\textbf{Repository:} \url{https://github.com/<user>/club-library}
\end{infobox}

\tableofcontents
\vspace{1em}

% ============================================================
\section{Introduction and motivation}
% ============================================================

Small clubs often manage their library with a spreadsheet: one row per book, a
column scribbled with the current borrower. This breaks down as soon as two
people edit the file, and it cannot answer questions such as \emph{``which books
are currently on loan?''} without manual counting.

\textbf{Club Library} replaces that spreadsheet with a proper relational system.
It manages \textbf{books}, \textbf{members} and \textbf{loans}, exposes the data
through a REST API and offers a small desktop frontend for the librarian.

\begin{beispielbox}
\textbf{Usage scenario.} A member returns two books. The librarian opens the
frontend, searches for the member, sees their two open loans and marks both as
returned. The stock counters update immediately, and the ``available books''
view now lists both titles again.
\end{beispielbox}

% ============================================================
\section{Data model}
% ============================================================

The domain has three entity types. A \textbf{member} can borrow many
\textbf{books} over time, and a book can be borrowed by many members -- an
\texttt{N:M} relationship that is resolved by the \textbf{loan} entity (L 02,
L 03).

\begin{beispielbox}
\textbf{ER model (textual).}
\begin{itemize}[leftmargin=1.4em, topsep=2pt]
  \item \textbf{member} (1) \;--\; borrows \;--\; (N) \textbf{loan}
  \item \textbf{book} (1) \;--\; is borrowed in \;--\; (N) \textbf{loan}
\end{itemize}
Thus \textbf{member} \texttt{N:M} \textbf{book}, resolved through
\textbf{loan}.
\end{beispielbox}

\subsection{Relational schema}

Derived from the ER model, with primary keys underlined and foreign keys in
italics:

\begin{itemize}[leftmargin=1.4em]
  \item \texttt{member}(\underline{id}, name, email)
  \item \texttt{book}(\underline{id}, title, author, isbn, stock)
  \item \texttt{loan}(\underline{id}, \emph{member\_id}, \emph{book\_id},
        borrowed\_on, due\_on, returned\_on)
\end{itemize}

\textbf{Normalisation.} Every non-key attribute depends only on the whole
primary key of its table, and there are no transitive dependencies: a member's
email lives only in \texttt{member}, a book's author only in \texttt{book}, and
a loan references members and books by their keys rather than duplicating their
data. The schema is therefore in \textbf{3NF} (L 04).

% ============================================================
\section{Database (PostgreSQL)}
% ============================================================

The schema and seed data are created by an init script that PostgreSQL runs on
first start (L 07, L 09).

\begin{lstlisting}[language=SQL, caption={init.sql (excerpt) -- schema and seed data}]
CREATE TABLE member (
    id    INTEGER GENERATED ALWAYS AS IDENTITY PRIMARY KEY,
    name  VARCHAR(120) NOT NULL,
    email VARCHAR(200) UNIQUE
);

CREATE TABLE book (
    id     INTEGER GENERATED ALWAYS AS IDENTITY PRIMARY KEY,
    title  VARCHAR(200) NOT NULL,
    author VARCHAR(120) NOT NULL,
    isbn   CHAR(13) UNIQUE,
    stock  INTEGER NOT NULL DEFAULT 1 CHECK (stock >= 0)
);

CREATE TABLE loan (
    id           INTEGER GENERATED ALWAYS AS IDENTITY PRIMARY KEY,
    member_id    INTEGER NOT NULL REFERENCES member(id),
    book_id      INTEGER NOT NULL REFERENCES book(id),
    borrowed_on  DATE NOT NULL DEFAULT CURRENT_DATE,
    due_on       DATE NOT NULL,
    returned_on  DATE
);
\end{lstlisting}

% ============================================================
\section{REST API (FastAPI)}
% ============================================================

The API is the only component that talks to the database (L 08). Read endpoints
are public; write endpoints require the header \texttt{X-API-Key} (cf.\ DBMS\_09).

\renewcommand{\arraystretch}{1.2}
\begin{tabularx}{\linewidth}{@{}p{1.6cm}p{4.4cm}Xc@{}}
\toprule
\textbf{Method} & \textbf{Path} & \textbf{Purpose} & \textbf{Key} \\
\midrule
\texttt{GET}  & \texttt{/books} & books + count of open loans (JOIN + agg.) & -- \\
\texttt{GET}  & \texttt{/members/\{id\}/loans} & open loans of a member (JOIN) & -- \\
\texttt{POST} & \texttt{/loans} & lend a book & yes \\
\texttt{POST} & \texttt{/returns} & return a book & yes \\
\bottomrule
\end{tabularx}

\medskip
The \texttt{/books} endpoint demonstrates a \textbf{JOIN} combined with an
\textbf{aggregation} (L 06): it counts how many copies of each book are
currently on loan.

\begin{lstlisting}[language=SQL, caption={Query behind GET /books (JOIN + aggregation)}]
SELECT b.id, b.title, b.author, b.stock,
       COUNT(l.id) FILTER (WHERE l.returned_on IS NULL) AS on_loan
FROM book b
LEFT JOIN loan l ON l.book_id = b.id
GROUP BY b.id
ORDER BY b.title;
\end{lstlisting}

The write path is guarded by a dependency that checks the API key -- adapted
from the DBMS\_09 exercise:

\begin{lstlisting}[language=Python, caption={API-key guard (adapted from DBMS\_09)}]
from fastapi import Header, HTTPException
import os

def require_api_key(x_api_key: str = Header(...)):
    if x_api_key != os.environ["API_KEY"]:
        raise HTTPException(status_code=401, detail="invalid API key")
\end{lstlisting}

% ============================================================
\section{Frontend (tkinter, packaged as .deb)}
% ============================================================

The frontend is a small tkinter application with two tabs: \textbf{Books} (a
searchable table backed by \texttt{GET /books}) and \textbf{Loans} (a form that
calls \texttt{POST /loans} and \texttt{POST /returns}). It never touches the
database directly -- it only calls the API (L 10).

At startup a \textbf{connection dialog} asks for the API URL and the
\texttt{X-API-Key}; the key is kept in memory for the session only and is never
written to disk.

\textbf{Packaging.} The application is frozen with PyInstaller and wrapped into a
Debian package with \texttt{fpm} (L 13), so the lecturer can install and run it
on the lecture server:

\begin{lstlisting}[caption={Building and installing the .deb}]
uv run pyinstaller --name club-library --onedir src/club_library/__main__.py
mkdir -p pkg/usr/bin && cp -r dist/club-library pkg/usr/bin/
fpm -s dir -t deb --name club-library --version 0.1.0 -C pkg .
sudo dpkg -i club-library_0.1.0_amd64.deb
\end{lstlisting}

% ============================================================
\section{Operation and deployment}
% ============================================================

The backend runs as two containers, orchestrated by Docker Compose and deployed
on the lecture server (L 09). Credentials live in a \texttt{.env} file that is
not committed.

\begin{lstlisting}[caption={docker-compose.yml (excerpt, after L 09)}]
services:
  postgres:
    image: postgres:16
    environment:
      POSTGRES_USER: ${POSTGRES_USER}
      POSTGRES_PASSWORD: ${POSTGRES_PASSWORD}
      POSTGRES_DB: ${POSTGRES_DB}
    volumes:
      - pg_data:/var/lib/postgresql/data
      - ./init.sql:/docker-entrypoint-initdb.d/init.sql
  api:
    build: ./api
    environment:
      API_KEY: ${API_KEY}
    env_file: .env
    ports: ["8000:8000"]
    depends_on: [postgres]
volumes:
  pg_data:
\end{lstlisting}

\textbf{Deployment steps on the lecture server.}
\begin{enumerate}[leftmargin=1.6em, topsep=2pt]
  \item \texttt{docker compose up -d} -- start database and API.
  \item \texttt{sudo dpkg -i club-library\_0.1.0\_amd64.deb} -- install the
        frontend.
  \item Launch \texttt{club-library}, enter the API URL and the
        \texttt{X-API-Key} in the connection dialog.
\end{enumerate}

% ============================================================
\section{Reflection}
% ============================================================

The separation into database, API and frontend paid off: the API-key guard and
the aggregation query could be tested with \texttt{curl} long before the
frontend existed. The trickiest part was \textbf{return handling} when several
copies of the same book are on loan at once -- solved by matching a return to
the oldest open loan of that book. Given more time, transaction isolation
(L 11) would make concurrent returns fully safe.

% ============================================================
\section*{Sources and reuse}
% ============================================================
\addcontentsline{toc}{section}{Sources and reuse}

Reused with attribution, as permitted:
\begin{itemize}[leftmargin=1.4em, topsep=2pt]
  \item API-key guard and Compose file structure -- adapted from
        \textbf{DBMS\_09} and \textbf{lecture 09}.
  \item tkinter connection dialog and \texttt{Treeview} layout -- adapted from
        \textbf{DBMS\_09 (frontend)} and \textbf{lecture 10}.
  \item PyInstaller / \texttt{fpm} packaging commands -- after \textbf{lecture
        13}.
\end{itemize}

\end{document}
